\section*{Introduction}
\addcontentsline{toc}{section}{Introduction}

Through a prescribed point of a singular variety, can one choose a Cartier hypersurface whose log canonical threshold is bounded away from zero? For a Fano type contraction $f:X\to Z$, the corresponding question asks for a hypersurface on $Z$ whose pullback has controlled singularities along the entire fibre. The hypersurface may depend on the contraction and the point. The desired lower bound depends only on the dimension of $X$ and a positive lower bound for its log discrepancies.

We work over an algebraically closed field of characteristic zero. A \emph{contraction} is a projective surjective morphism of normal integral varieties with connected fibres, expressed by $f_*\mathcal O_X=\mathcal O_Z$. It is of \emph{Fano type} if an effective rational boundary $C$ exists such that $(X,C)$ is klt and $-(K_X+C)$ is ample over $Z$. This auxiliary boundary need not equal the boundary in the following theorem.

A rational pair $(X,B)$ has effective rational boundary $B$ and rationally Cartier log canonical divisor $K_X+B$. If $p:W\to X$ is a smooth birational model and $E$ is a prime divisor on $W$, our discrepancy convention is
\[
a(E;X,B)=1+\operatorname{coeff}_E\bigl(K_W-p^*(K_X+B)\bigr).
\]
Thus $\epsilon$-lc means $a(E;X,B)\geq\epsilon$ for every such $E$; lc means that they are nonnegative; and klt means that they are strictly positive.

The obstruction is visible in a single ratio. If $D$ has local equation $g$ on the base, log canonicity after adding $t f^*D$ requires
\[
 t\le\frac{a(E;X,B)}{\ord_E(f^*g)}
\]
for every prime divisor $E$ over the relevant inverse image with positive denominator. A discrepancy bound alone does not control those orders. Even on a smooth curve, the divisor $m[z]$ has threshold $1/m$. The problem is therefore to choose an equation through $z$ for which all these inequalities hold with one uniform positive $t$.

\begin{samepage}
\begin{maintheorem}[Uniform Cartier sections]\label{thm:main}
For every integer $d>0$ and every real number $\epsilon>0$, there is a real number $\tau=\tau(d,\epsilon)>0$ with the following property. Let $(X,B)$ be an $\epsilon$-lc rational pair of dimension $d$, and let $f:X\to Z$ be a Fano type contraction of normal integral quasi-projective varieties with $\dim Z>0$. Suppose that $-(K_X+B)$ is nef over $Z$. For every closed point $z\in Z$, there are an open neighbourhood $U$ of $z$ and a nonzero effective Cartier divisor $D$ on $U$, with $z\in\Supp D$, such that
\[
\bigl(f^{-1}(U),B|_{f^{-1}(U)}+\tau f^*D\bigr)
\quad\text{is log canonical.}
\]
\end{maintheorem}
\end{samepage}

The conclusion of Theorem~\ref{thm:main} concerns the whole inverse image of a neighbourhood of $z$. The constant is independent of Cartier indices, boundary denominators, and any choice of polarization. It is non-effective; the divisor and the neighbourhood may depend on the contraction and the point.

\paragraph{Cartier sections and complements.}
For the identity contraction, the question above concerns arbitrary $\epsilon$-lc singularities. Equivalently, it asks for a uniform positive lower bound for the log canonical threshold of the maximal ideal at a closed point; the passage to a general element of that ideal is recalled in Lemma~\ref{lem:H-field-reduction}. The general contraction requires the corresponding bound for the pulled-back point ideal. This is a different local question from bounding the discriminant coefficients in the canonical bundle formula: those coefficients are defined by thresholds over generic points of divisors on birational models of the base. A Cartier section through a prescribed point must control the whole fibre at once.

The complement formulation asks for a uniformly bounded complement index. Over a neighbourhood of a base point, a \emph{monotonic klt $n$-complement} of $(X,B)$ is an effective rational divisor $\Lambda\ge B$ such that $(X,\Lambda)$ is klt and
\[
 n(K_X+\Lambda)\sim 0.
\]
Here linear equivalence includes the assertion that the displayed divisor is Cartier. In Shokurov's boundedness conjecture, as formulated in \cite[Conjecture~1.1]{Chen2025}, the index $n$ depends only on the total dimension $d$, the discrepancy bound $\epsilon>0$, and a fixed finite set $\mathfrak R\subset[0,1]\cap\Q$ containing the coefficients of $B$. The geometric hypotheses are those of Theorem~\ref{thm:main}, with the zero-dimensional base also allowed.

Complements entered birational geometry through Shokurov's work on threefold log flips, including bounded local lc complements for surfaces \cite[Section~5]{ShokurovFlips}. His surface theorem and bounded-complement conjectures developed the relation between mild singularities and bounded complement indices \cite[Conjecture~1.3 and Main Theorem~1.4]{ShokurovSurfaces}. The stronger expectation that bounded complements preserve a positive lower bound on log discrepancies appears in his 2004 formulation \cite{ShokurovBounded2004}. Birkar proved boundedness of lc complements in arbitrary dimension for Fano type pairs with hyperstandard coefficients \cite[Theorems~1.7--1.8]{BirkarComplements}. Those lc complements provide the bounded initial index used below. A klt complement must have strictly positive discrepancies, and obtaining this additional property uniformly is the issue addressed here.

Several earlier results describe the reach of the Cartier-section problem. Sankaran and Santos bounded the smallest weight of an $\epsilon$-lc weighted blowup of affine space \cite[Theorem~1.3]{SankaranSantos}, proving the weighted-blowup case of Birkar's question. For torus-invariant boundaries with coefficients in a fixed DCC set (one with no infinite strictly decreasing sequence), Ambro constructed invariant Cartier sections through torus-fixed base points in the toric setting \cite[Theorem~3.12, arXiv version~2]{AmbroToric}. Under his hyperstandard coefficient hypothesis, his Theorem~1.2 also gives bounded complements preserving a positive discrepancy bound. Birkar's control of base singularities gives the Cartier assertion over curves, together with bounded klt complements in the curve-base case \cite[Theorem~1.1 and Corollary~1.4]{BirkarCY}. These curve and toric results illustrate two distinct ways to control the orders appearing in the threshold inequalities.

Over $\C$, Chen proves the decisive general reduction. For each bound $\ell$ on the base dimension, bounded klt complements, the Cartier-section assertion, and its identity-contraction case are equivalent; on the complement side, birational contractions with zero boundary suffice \cite[Theorem~1.6]{Chen2025}. He proves the surface local assertion with coefficient $\epsilon^2/24$, and hence both conjectures for bases of dimension at most two \cite[Theorem~1.7 and Corollary~1.8]{Chen2025}. He also removes the coefficient restriction in the toric Cartier-section assertion \cite[Theorem~1.11]{Chen2025}. We use the precise local formulation of the Birkar--Shokurov conjecture in \cite[Conjecture~1.3]{Chen2025}; the relatively ample formulation also appears in \cite[Conjecture~1.27]{BirkarSurvey}.

Theorem~\ref{thm:main} proves the rational-boundary assertion in characteristic zero, and Corollary~\ref{cor:real-boundaries} below gives the real-boundary formulation over $\C$. The conjecture is therefore resolved positively in these settings. The new input to Chen's reduction is the following birational theorem.

\begin{maintheorem}[Birational klt complements]\label{thm:birational}
For each $d>0$ and $\epsilon>0$, there is a positive integer $n=n(d,\epsilon)$ such that the following holds. Let $X$ be an $\epsilon$-lc rationally Gorenstein variety of dimension $d$, and let $f:X\to Z$ be a birational Fano type contraction with $-K_X$ nef over $Z$. For every closed $z\in Z$, there is an effective rational divisor $\Delta$ over a neighbourhood of $z$ such that $(X,\Delta)$ is klt there and $n(K_X+\Delta)$ is Cartier and linearly trivial there, after shrinking the neighbourhood.
\end{maintheorem}

The companion article \cite[Theorem~1.1 and Corollary~1.2]{OpenAIComplements}
treats bounded klt complements for Fano contractions and their
finite-rational-coefficient extension by an independent argument.
The present proof establishes Theorem~\ref{thm:birational} directly,
through its own character and pinning obstruction.

Applying Chen's Theorem~1.6 in arXiv version~4, implication (2)$\Rightarrow$(3), transfers Theorem~\ref{thm:birational} to Theorem~\ref{thm:main}. For a fixed total dimension $d$, use the base-dimension bound $\ell=d$ and the birational theorem in every dimension at most $d$, for every positive discrepancy gap. Implication (2)$\Rightarrow$(1) of the same theorem also supplies the finite-rational-coefficient complement assertion stated above over $\C$. Its Cartier-section conclusion already allows real boundaries, giving the following extension.

\begin{maincorollary}[Real boundaries]\label{cor:real-boundaries}
Over $\C$, the conclusion of Theorem~\ref{thm:main} holds for effective real boundaries $B$ with $K_X+B$ real Cartier, with the same uniform dependence on $d$ and $\epsilon$.
\end{maincorollary}

The reduction requires rationalization only in the zero-boundary birational case: an integral Weil divisor that is real Cartier is rationally Cartier, and a real klt log Fano witness can be chosen rational. Section~\ref{sec:H} supplies these details, the reduction from arbitrary algebraically closed characteristic-zero fields to $\C$ for Theorem~\ref{thm:main}, and the passage between point-ideal thresholds and individual Cartier divisors. No approximation of the arbitrary nef real boundary in Corollary~\ref{cor:real-boundaries} is required.

\paragraph{The argument.}
Suppose that Theorem~\ref{thm:birational} fails. Section~\ref{sec:H} uses bounded lc complements to construct a sequence of complex klt germs $x_i\in S_i$ of bounded dimension, with finite cyclic groups $\Gamma_i$ fixing $x_i$. The quotient maps
\[
 S_i\longrightarrow S_i/\Gamma_i
\]
are quasi-\'etale, meaning \'etale outside subsets of codimension at least two, and the quotient germs have a fixed positive discrepancy gap for primitive prime orders. On $S_i$, a nonzero regular function $h_i$ and a nowhere-vanishing canonical generator have the same character $\chi_i$. The pair $(S_i,\divisor(h_i))$ is lc, and a centred quasi-monomial valuation $T_i$ satisfies
\[
 A_{S_i}(T_i)=T_i(h_i)=1,
\]
where $A_{S_i}$ is the log-discrepancy function with zero boundary. Failure of bounded klt complements gives the stronger condition
\[
 \forall m\in\Z_{>0}\ \exists i(m)\ \forall i\ge i(m):\qquad
 T_i(s)\ge m
 \quad\text{for every nonzero regular $s$ of character $\chi_i^m$.}
\]
Thus the order bound holds for every section in each fixed character degree, rather than for a selected sequence of sections. The reduction uses a canonical cover in the isomorphism case and a signed anti-canonical grading for other birational contractions.

Character relations become concrete after division by powers of $h_i$. For a section $s$ of character $\chi_i^m$, the rational function $Y=s/h_i^m$ is $\Gamma_i$-invariant and
\[
 T_i(Y)=T_i(s)-m.
\]
Requiring nonnegative order for both $Y$ and $Y^{-1}$ forces the exact equality $T_i(Y)=0$. The construction uses pairs of ratios with opposite orders to impose these two inequalities. The induction builds and retains relations of this kind. It optimizes valuations on the constrained $h_i$-lc slice, passes to finitely generated central algebras, and imposes further invariant relations while preserving the algebraic independence of the variables already obtained. Either it produces another independent variable or the remaining character direction yields a quotient prime of arbitrarily small discrepancy. Dimension bounds the first process, while the discrepancy gap excludes the second. Theorem~\ref{thm:G-pinning} formalizes this obstruction, and Section~\ref{sec:H} uses it to prove the birational theorem before applying Chen's equivalence.

The proof requires uniform geometric estimates and exact control of the valuation constraints. Sections~\mbox{\ref{sec:A}--\ref{sec:C}} establish the discrepancy comparisons, phase nonvanishing, and saturation statements for fields obtained by cutting off small discrepancy directions. Their geometric inputs include Birkar's boundedness of weak log Fano varieties \cite[Theorem~1.1]{BirkarBAB}, boundedness of rationally connected generalized Calabi--Yau varieties up to isomorphism in codimension one \cite[Theorem~1.7]{BirkarCY}, and effective birationality for nef and big integral Weil divisors \cite[Theorem~1.1]{BirkarPolarised}. The latter theorem also requires the polarization minus the canonical divisor to be pseudo-effective. These results apply without a prior bound on the nef divisors' Cartier indices. Filipazzi's generalized canonical bundle formula supplies the generalized-pair structures on successive bases \cite[Theorem~1.4]{Filipazzi}; the uniform discrepancy comparisons needed here are established in Section~\ref{sec:A}.

For the local optimization, the normalized-volume theory of Li, Blum, Xu, Li--Xu, and Xu--Zhuang supplies existence, quasi-monomiality, uniqueness, and stable degeneration for ordinary klt-pair minimizers \cite{LiIzumi,Blum,XuQuasimonomial,LiXuStable,XuZhuangUnique,XuZhuangStable}. The order inequalities above impose additional constraints. Section~\ref{sec:E} constructs an effective rational klt pair whose ordinary minimizer lies exactly on the constrained slice; only then is the stable-degeneration theorem applied. Section~\ref{sec:F} transfers the resulting equalities and signs through later degenerations and changes of grading. Together with the character counts of Section~\ref{sec:D}, these exact transfers ensure that each step of the induction in Section~\ref{sec:G} retains the earlier relations.

\paragraph{Conventions on generalized pairs and sequences.}
For a generalized pair, a rationally Cartier nef b-divisor $\mathbf M$ is determined and nef on a model $\pi:Y\to X$. On that model the boundary is defined by
\[
K_Y+B_Y+\mathbf M_Y=\pi^*(K_X+B+\mathbf M_X).
\]
Discrepancies on higher models are the ordinary discrepancies of the resulting, possibly noneffective, boundary. The order of a rationally Cartier b-divisor is computed after pulling it back from a determining model. Unless expressly stated otherwise, rational divisor data are used for MMPs.

Orders and discrepancies extend homogeneously to positive multiples of divisorial valuations. A lower discrepancy bound on genuine singularities is always imposed on primitive orders $\ord_E$, not on all of their positive rescalings. This distinction is essential when valuations on covers, quotients, and subfields are compared.

The auxiliary arguments concern sequences in bounded dimension. The symbols $O(1)$ and $o(1)$ refer to those sequences, with additional fixed parameters indicated at the point of use. Sequence conditions are required on every subsequence. We pass to further subsequences when proving a contradiction and choose auxiliary growth rates slowly enough to satisfy the preceding finite sets of requirements. Each uniform conclusion keeps this order of choices explicit. In Sections~\ref{sec:E}--\ref{sec:H}, germs are complex algebraic germs. An $h$-lc valuation means a valuation of zero discrepancy for the pair with boundary $\divisor(h)$, using a power to descend the rational divisor through a finite quotient when necessary.
